\documentclass{standalone}
\usepackage{circuitikz}
\usetikzlibrary{calc}
\definecolor{circuitnotemark}{HTML}{FE00FE}
\begin{document}
\begin{circuitikz}[american, line width=0.8pt]
\ctikzset{bipoles/length=1.2cm}
\ctikzset{grounds/scale=1.36}
\coordinate (x1y2) at (0,-2);
\coordinate (x1y5) at (0,-8);
\coordinate (x2y1) at (2,0);
\coordinate (x4y2) at (6,-2);
\coordinate (x7y2) at (12,-2);
\coordinate (x7y5) at (12,-8);
\coordinate (x7y1) at (12,0);
\coordinate (x2y2) at (2,-2);
\draw (x1y2) to[esource, l_=$\mathrm{W1}$] (x1y5); % line 3
\draw (-0.18,-5) -- ++(0,0.18) -- ++(0.18,0) -- ++(0,-0.36) -- ++(0.18,0) -- ++(0,0.18); % line 3
\node[ground] at (x1y5) {}; % line 4
\draw (x2y1) node[ocirc]{} node[above left]{CH1}; % line 5
\draw (x4y2) to[R, l_=$R_{1}$, a^=$1\,\mathrm{k}\Omega$] (x7y2); % line 6
\draw (x7y2) to[C, l_=$C_{1}$, a^=$1\,\mu\mathrm{F}$] (x7y5); % line 7
\node[ground] at (x7y5) {}; % line 8
\draw (x7y1) node[ocirc]{} node[above left]{CH2}; % line 9
\draw (x1y2) -- (x2y2); % line 11
\draw (x2y2) -- (x4y2); % line 12
\draw (x2y1) -- (x2y2); % line 13
\draw (x7y1) -- (x7y2); % line 14
\node[circ] at (x7y2) {};
\node[circ] at (x2y2) {};
\node[anchor=south west, inner sep=2pt, yshift=4pt, circuitnotemark, font=\LARGE\bfseries] (circuittitle) at (current bounding box.north west) {X};
\path (circuittitle.south west) rectangle ++(7.315,0.853);
\node[anchor=north east, inner sep=2pt, circuitnotemark, font=\large] (circuitstamp) at ([yshift=-0.593cm]current bounding box.south east) {X};
\path (circuitstamp.north east) rectangle ++(-5.08,-0.593);
\end{circuitikz}
\end{document}
